\section{Experimental Linear-in-Means Model with Subject-Specific Instruments}

Suppose individuals are organized into groups of size \( n \) (e.g., classrooms) within schools. For each individual \( i \) we observe an outcome \( y_i \) (generated from a latent variable), an exogenous regressor \( x_i \), and a subject label \( s_i \) (e.g., Math, English, Physics, \(\dots\), from a set \(\mathcal{S}\) of 20 subjects). Let \( \mu_{s_i} \) denote the subject-specific fixed effect. The baseline latent variable model is
\[
y_i^* = \alpha + \beta\, x_i + \mu_{s_i} + \lambda_{\text{eff}} \, (Wy)_i + \varepsilon_i,
\]
where
\[
\lambda_{\text{eff}} = \lambda_C + \delta \, D_i,
\]
with \( D_i \) a (group-level) treatment indicator and
\[
(Wy)_i = \sum_{j\neq i} w_{ij}\, y_j, \quad w_{ij} = \frac{1}{n-1}, \quad w_{ii}=0.
\]
The observed binary outcome is given by
\[
y_i = \mathbf{1}\{ y_i^* > 0 \}.
\]

Pooling individuals, the structural model to be estimated as a linear probability model (LPM) is
\[
y_i = \alpha + \beta\, x_i + \lambda_C\,(Wy)_i + \delta\, \bigl[ D_i\,(Wy)_i \bigr] + \eta_i,
\]
where \(\eta_i\) is an error term.

\subsection*{Endogeneity and IV Identification}

The peer regressor \((Wy)_i\) is endogenous because it is determined jointly with \( y_i \) (the reflection problem). In our baseline IV strategy we might use the leave--one--out average of \(x\),
\[
z_i = (Wx)_i = \sum_{j \neq i} w_{ij}\, x_j,
\]
and its interaction \(D_i\, z_i\) as instruments. However, we wish to refine the instrument by exploiting the fact that subject content influences the outcome. 

Define the set of peers within the same group who share the same subject as
\[
\mathcal{N}_i(s) = \{ j \neq i : s_j = s_i \}.
\]
We then define the subject-specific leave--one--out average of \( x \) as:
\[
z_i^{(s)} = \begin{cases}
\displaystyle \frac{1}{|\mathcal{N}_i(s)|} \sum_{j\in \mathcal{N}_i(s)} x_j, & \text{if } |\mathcal{N}_i(s)| > 0, \\[1mm]
x_i, & \text{if } |\mathcal{N}_i(s)| = 0.
\end{cases}
\]
We use \( z_i^{(s)} \) and \( D_i\, z_i^{(s)} \) as instruments for \((Wy)_i\) and \(D_i\,(Wy)_i\), respectively. Thus, the instrument vector for individual \(i\) is:
\[
Z_i = \begin{pmatrix} 1 \\ x_i \\ z_i^{(s)} \\ D_i\, z_i^{(s)} \end{pmatrix}.
\]

\subsection*{Two-Stage Least Squares (2SLS) Estimation}

The IV strategy is implemented in two stages:
\begin{enumerate}
    \item \textbf{First Stage:} Regress the endogenous variables
    \[
    \begin{pmatrix} (Wy)_i \\ D_i\,(Wy)_i \end{pmatrix}
    \]
    on the instrument vector \( Z_i \) to obtain the predicted values
    \[
    \begin{pmatrix} \widehat{(Wy)}_i \\ \widehat{D_i\,(Wy)}_i \end{pmatrix}.
    \]
    
    \item \textbf{Second Stage:} Regress \( y_i \) on an intercept, \( x_i \), \( \widehat{(Wy)}_i \), and \( \widehat{D_i\,(Wy)}_i \):
    \[
    y_i = \alpha + \beta\, x_i + \lambda_C\, \widehat{(Wy)}_i + \delta\, \widehat{D_i\,(Wy)}_i + \eta_i.
    \]
    The coefficients on \( \widehat{(Wy)}_i \) and \( \widehat{D_i\,(Wy)}_i \) provide consistent estimates of \( \lambda_C \) and \( \delta \), respectively.
\end{enumerate}

\subsection*{Summary}

To summarize, our identification strategy for the experimental linear-in-means model with subject heterogeneity is as follows:
\begin{align*}
y_i^* &= \alpha + \beta\, x_i + \mu_{s_i} + \left[\lambda_C + \delta\, D_i\right] (Wy)_i + \varepsilon_i, \\
y_i &= \mathbf{1}\{ y_i^* > 0 \}, \\
(Wy)_i &= \sum_{j \neq i} w_{ij}\, y_j, \quad w_{ij}=\frac{1}{n-1}, \; w_{ii}=0, \\
z_i^{(s)} &= \frac{1}{|\mathcal{N}_i(s)|} \sum_{j\in\mathcal{N}_i(s)} x_j, \quad \text{(with the convention that if } |\mathcal{N}_i(s)| = 0, \text{ then } z_i^{(s)} = x_i\text{)}, \\
Z_i &= \begin{pmatrix} 1 \\ x_i \\ z_i^{(s)} \\ D_i\, z_i^{(s)} \end{pmatrix}.
\end{align*}
The structural model for estimation is
\[
y_i = \alpha + \beta\, x_i + \lambda_C\,(Wy)_i + \delta\, \left[ D_i\,(Wy)_i \right] + \eta_i,
\]
and we use 2SLS with \(Z_i\) as instruments for \((Wy)_i\) and \(D_i\,(Wy)_i\).
